\section{The effective potential with the self-interaction}
\label{sec:cjtcontact}

Section~\ref{sec:cjt} applied the effective action of Cornwall, Jackiw and Tomboulis to the model
in which the Vorton-Scalar alone provides mass and interaction. This section applies the same
method to the self-interaction of the fermion field: first without the Vorton-Scalar, where it must
reproduce the Kohn-Sham results of Sec.~\ref{sec:ks} and goes beyond them by counting the filled
states of negative energy; then with the bare Dirac mass, the self-interaction and the Vorton-Scalar
together.

\subsection{The effective potential of uniform matter}

The CJT effective action is a functional of the full propagator $S$ of the fermions,
\begin{equation}
\Gamma[S]=-i\,\mathrm{Tr}\ln S^{-1}-i\,\mathrm{Tr}\big(S_0^{-1}S-1\big)+\Gamma_2[S],
\label{eq:GammaCJT}
\end{equation}
where $S_0$ is the propagator of free particles of the mass $m$ and $\Gamma_2$ is the sum of the
vacuum diagrams that cannot be cut in two by cutting two lines, drawn with the full propagator. For
the contact interaction $\Lself$ the lowest order of $\Gamma_2$ consists of two diagrams, each with
two loops that meet at one vertex: the direct (Hartree) and the exchange (Fock) diagram. At this
order the stationarity condition $\delta\Gamma/\delta S=0$ says that the self-energy
$S_0^{-1}-S^{-1}$ does not depend on the momentum. For uniform matter at rest it is a combination of
the unit matrix and of $\beta$ (Sec.~\ref{sec:hf}): the full propagator is that of free particles of
a mass $\mst$ whose energies are shifted by a constant $V_0$. The effective action per unit volume
and time then becomes an ordinary function of the two numbers $\mst$ and $V_0$, the \emph{effective
potential}. At zero temperature and at the chemical potential $\mu$,
\begin{align}
\Omega_{\mathrm{CJT}}&(\mst,V_0;\mu)=-\frac{(\mst-m)^{2}}{2\,\GS}-\frac{V_0^{2}}{2\,\GV}\nonumber\\
&-\nu\!\int\!\frac{d^{D}k}{(2\pi)^{D}}\,\big(\mu-V_0-E^{*}\big)\,\theta\big(\mu-V_0-E^{*}\big),
\label{eq:OmegaCJT}
\end{align}
with $E^{*}=(k^{2}+\mstq)^{1/2}$, the step function $\theta$, and the \SIG{scalar and vector coupling
$\GS$, $\GV$ of matter of one kind of particle with $\nu$ internal states (Sec.~\ref{sec:ksmethod})}{couplings
$\GS$, $\GV$ of Eq.~(\ref{eq:GHF})}. The integral is the grand potential of free particles of the mass $\mst$ at
the chemical potential $\mu-V_0$; the two other terms remove the interaction energy that the
integral counts twice. The two stationarity conditions are
\begin{equation}
\begin{aligned}
\frac{\partial\Omega_{\mathrm{CJT}}}{\partial V_0}&=0:\quad V_0=\GV\,n,\\
\frac{\partial\Omega_{\mathrm{CJT}}}{\partial\mst}&=0:\quad\mst=m+\GS\,\nS ,
\end{aligned}
\label{eq:CJTstat}
\end{equation}
the two equations (\ref{eq:hKS}) of the Kohn-Sham method. At the stationary point
$-\Omega_{\mathrm{CJT}}$ is the pressure and $\Omega_{\mathrm{CJT}}+\mu\,n$ the energy density.
\emph{The Kohn-Sham method and the CJT effective potential at this order are two formulations of the
same approximation}: the first works at a fixed density and solves the gap equation, the second
works at a fixed chemical potential and searches the stationary point of a function.

The program evaluates the two with separate code (Gauss-Legendre integrals and a root search in the
one, Simpson integrals and a minimization in the other) for {{GEN:cjt_n_states}} states in $D=7$,
$4$ and $3$ dimensions (Table~\ref{tab:cjtks}). Mass, density, energy density and pressure agree to
{{GEN:cjt_ks_dev}} (sum of the four relative deviations); the independent script repeats the
comparison with its own integrals and finds the same.

\begin{table}
\caption{\label{tab:cjtks}The Kohn-Sham method (KS) and the CJT effective potential compared, without
the filled states of negative energy (units $m=1$): the mass of the particles in matter by both
methods, and $-\Omega_{\mathrm{CJT}}$ at its stationary point divided by the Kohn-Sham pressure $p$,
for ten of the {{GEN:cjt_n_states}} states computed (the data files hold eleven digits).}
\begin{ruledtabular}
\begin{tabular}{ccccccc}
$D$ & $\GS$ & $\GV$ & $\kF$ & $\mst$ (KS) & $\mst$ (CJT) & $-\Omega_{\mathrm{CJT}}/p$ \\
\hline
{{ROWS:tbl_cjt_vs_ks}}
\end{tabular}
\end{ruledtabular}
\end{table}

The effective potential says more than the gap equation in one respect. At a given chemical
potential the stationary point with the lowest $\Omega_{\mathrm{CJT}}$, that is, with the highest
pressure, is the state of equilibrium (the vector field $V_0$ is taken at its stationary value, a
maximum; the minimum is over $\mst$). Of the {{GEN:cjt_n_positive}} states of the comparison with a
positive pressure, {{GEN:cjt_n_global}} are that minimum. In the others a denser state has the higher
pressure at the same chemical potential: the Kohn-Sham state of that density is metastable, a dilute
gas that can lower its energy by condensing into saturated matter.
%%<1644>

\emph{The two halves of the field.} For the matter of Sec.~\ref{sec:ksreduced}, whose two halves have
different masses and energy shifts, the effective potential is a function of four numbers,
\begin{equation}
\begin{aligned}
\Omega_{\mathrm{CJT}}={}&\sum_\pm\Omega_{\mathrm{half}}\big(\mstpm,\mu-V_{0\pm}\big)\\
&-\frac{\big[\tfrac12(m^{*}_++m^{*}_-)-m\big]^{2}}{2\,\GStot}-\frac{\big[\tfrac12(m^{*}_+-m^{*}_-)\big]^{2}}{2\,\GSdif}\\
&-\frac{\big[\tfrac12(V_{0+}+V_{0-})\big]^{2}}{2\,\GVtot}-\frac{\big[\tfrac12(V_{0+}-V_{0-})\big]^{2}}{2\,\GVdif},
\end{aligned}
\label{eq:OmegaHalves}
\end{equation}
where $\Omega_{\mathrm{half}}(\mst,\mu)=-4\int d^{D}k\,(\mu-E^{*})\,\theta(\mu-E^{*})/(2\pi)^{D}$ is the
grand potential of free particles of one half, with four internal states. Its four stationarity
conditions are the equations (\ref{eq:hKShalves}) with the densities (\ref{eq:denshalves}). For
{{GEN:fm_cjt_n_states}} states of both forms ({{GEN:fm_cjt_n_equal}} with equal halves,
{{GEN:fm_cjt_n_unequal}} with unequal halves, {{GEN:fm_cjt_n_one}} with one half filled) the
stationary point, found by Newton's method from displaced start values, agrees with the Kohn-Sham
state in both masses, density, energy density and pressure to {{GEN:fm_cjt_dev}}: the two methods
are equivalent in both forms.
%%</1644>

\subsection{The filled states of negative energy}
\label{sec:sea}

The Kohn-Sham calculation counts the particles of matter only. The term
$-i\,\mathrm{Tr}\ln S^{-1}$ of Eq.~(\ref{eq:GammaCJT}) also contains the filled states of negative
energy of the vacuum, whose energy depends on the mass of the particles. Because the contact
interaction is an effective interaction, their momenta are counted up to a cutoff $\Lambda$:
\begin{equation}
\begin{aligned}
\Omega_{\mathrm{CJT}}\to{}&\Omega_{\mathrm{CJT}}-\nu\!\int_{|\vec k|<\Lambda}\!\frac{d^{D}k}{(2\pi)^{D}}\\
&\times\Big[\sqrt{k^{2}+\mstq}-\sqrt{k^{2}+m^{2}}\Big].
\end{aligned}
\label{eq:sea}
\end{equation}
The filled states contribute the scalar density
$-\nu\int_{|\vec k|<\Lambda}d^{D}k\,(\mst/E^{*})/(2\pi)^{D}$ to the gap equation. The mass parameter
of the Lagrangian is then no longer the mass of a particle in vacuum and is called the \emph{bare
mass} $m_{\mathrm{bare}}$; the mass in vacuum, $m_{\mathrm{vac}}$, is the solution of
\begin{equation}
\begin{aligned}
m_{\mathrm{vac}}&=m_{\mathrm{bare}}+\GS\,\langle\bar\psi\psi\rangle_{\mathrm{vac}},\\
\langle\bar\psi\psi\rangle_{\mathrm{vac}}&=-\nu\!\int_{|\vec k|<\Lambda}\!\frac{d^{D}k}{(2\pi)^{D}}\,\frac{m_{\mathrm{vac}}}{\sqrt{k^{2}+m_{\mathrm{vac}}^{2}}} .
\end{aligned}
\label{eq:mvac}
\end{equation}
The scalar density of the vacuum is negative, so that an attraction ($\GS<0$) \emph{raises} the mass
above the bare mass.

\emph{A critical coupling.} For $m_{\mathrm{bare}}=0$ Eq.~(\ref{eq:mvac}) always has the solution
$m_{\mathrm{vac}}=0$. A second solution, of lower energy, exists if the attraction exceeds a
critical value, $\GS<G_c$. Expanding the integral for small $m_{\mathrm{vac}}$, with the area
$A_D=D\,V_D$ of the unit sphere,
\begin{equation}
\begin{aligned}
G_c&=-\frac{(2\pi)^{D}\,(D-1)}{\nu\,A_D\,\Lambda^{D-1}},\\
G_c\,\Lambda^{D-1}&=-90\pi^{4},\ -3\pi^{2},\ -\frac{\pi^{2}}{2}\quad\text{for }D=7,\ 4,\ 3 .
\end{aligned}
\label{eq:Gc}
\end{equation}
\emph{A field without a mass term acquires a mass by its own self-interaction when the scalar
attraction exceeds $G_c$.} This is the mechanism of Nambu and Jona-Lasinio \cite{NJL1961};
the minimization of the effective potential gives no mass two per cent below $G_c$ and a mass two
per cent above it in all three cases.

\emph{On the deflating spacetime the critical coupling is crossed at a definite moment.} In the
sector of fields that do not depend on $x5,x6,x7$ the coupling grows as $e^{\td}$ while the cutoff,
a momentum along the flat coordinates, does not change. A field that is massless at early times
therefore acquires a mass at the clock reading $\td=\ln[G_c/\GS(x4_{\mathrm{ref}})]$, and from then
on the mass grows [Fig.~\ref{fig:cjtsea}(a), Table~\ref{tab:cjtvac}]. With a small bare mass the
transition is smooth. The mass reaches the cutoff about half a unit of the clock after the critical
coupling is crossed; beyond that point the effective theory says nothing.\SIG{ (Table~\ref{tab:cjtvac}
and Fig.~\ref{fig:cjtsea} are for one kind of particle with eight internal states, that is, for form Q
of the $4+4$ model and for the $7+1$ model; the two halves of the field in form C follow below.)}{}

\begin{table}
\caption{\label{tab:cjtvac}The mass of the fermions in vacuum, $m_{\mathrm{vac}}/\Lambda$, with the
filled states of negative energy, along the deflation clock counted from the moment at which
$\GS=G_c$, for a field without bare mass and for the bare mass $0.05\,\Lambda$.}
\begin{ruledtabular}
\begin{tabular}{cccccc}
$D$ & $m_{\mathrm{bare}}/\Lambda$ & $\td=-1$ & $\td=0$ & $\td=0.24$ & $\td=0.48$ \\
\hline
{{ROWS:tbl_cjt_vacuum}}
\end{tabular}
\end{ruledtabular}
\end{table}

\begin{figure*}
\includegraphics[width=0.86\textwidth]{fig_cjt_sea}
\caption{\label{fig:cjtsea}The filled states of negative energy. (a) The mass of the fermions in
vacuum, in units of the cutoff, along the deflation clock counted from the moment at which the
scalar attraction equals the critical coupling, for $D=7$, $4$ and $3$ dimensions, without bare mass
and with the bare mass $0.05\,\Lambda$. (b) Energy per particle of matter relative to the vacuum, for
\SIG{an example on the $7+1$ spacetime}{the reference example of Sec.~\ref{sec:ksflat}} ($D=7$,
$\GS=-10^{5}$, $\GV=+10^{4}$ in units of the mass in vacuum) without the filled states and with them
for three cutoffs; the bare mass is chosen so that the mass in vacuum is the same.}
\end{figure*}

%%<1671>
\emph{The reference example with the filled states.} The coupling $\GS=-10^{5}\,m^{-6}$ of the
reference example of Sec.~\ref{sec:ksflat} exceeds the critical coupling $90\pi^{4}/\Lambda^{6}$ as
soon as the cutoff is larger than $\Lambda={{GENM:cjt_cut_limit}}\,m$: \emph{the calculation without
the filled states presupposes a cutoff below this value}, which is smaller than the Fermi momentum
of the saturated matter. With the filled states included, and with the bare mass adjusted so that
the mass in vacuum stays the same, the saturation point moves little for the two smaller cutoffs of
Table~\ref{tab:cjtsea}, and the binding grows for the largest [Fig.~\ref{fig:cjtsea}(b)]. The example is therefore meaningful only as an
effective theory with a low cutoff, and its numbers depend on the cutoff at the level that the table
shows.

\begin{table}
\caption{\label{tab:cjtsea}The reference example ($D=7$, $\GS=-10^{5}$, $\GV=+10^{4}$, units
$m_{\mathrm{vac}}=1$) with the filled states of negative energy below the cutoff $\Lambda$: the bare
mass that keeps the mass in vacuum, and the saturation point (Fermi momentum, mass of the particles
in matter, energy per particle relative to the vacuum).}
\begin{ruledtabular}
\begin{tabular}{ccccc}
$\Lambda$ & $m_{\mathrm{bare}}$ & $\kF$ & $\mst$ & $\varepsilon/n-m_{\mathrm{vac}}$ \\
\hline
{{ROWS:tbl_cjt_sea}}
\end{tabular}
\end{ruledtabular}
\end{table}
%%</1671>
%%<1644>
\emph{The vacuum of the two halves.} For the two halves of the field the gap equations of the
vacuum are
\begin{equation}
\begin{aligned}
\tfrac12(m^{*}_++m^{*}_-)-m_{\mathrm{bare}}&=\GStot\,\big(n_{S+}^{\mathrm{vac}}+n_{S-}^{\mathrm{vac}}\big),\\
\tfrac12(m^{*}_+-m^{*}_-)&=\GSdif\,\big(n_{S+}^{\mathrm{vac}}-n_{S-}^{\mathrm{vac}}\big),\\
n_{S\pm}^{\mathrm{vac}}&=-4\!\int_{|\vec k|<\Lambda}\!\frac{d^{D}k}{(2\pi)^{D}}\,\frac{\mstpm}{E^{*}_\pm},
\end{aligned}
\label{eq:vachalves}
\end{equation}
with the scalar densities $n_{S\pm}^{\mathrm{vac}}$ of the filled states of the two halves; where they
have several solutions, the vacuum is the one of lowest energy density
\begin{equation}
\begin{aligned}
\Omega_{\mathrm{vac}}={}&-4\sum_\pm\!\int_{|\vec k|<\Lambda}\!\frac{d^{D}k}{(2\pi)^{D}}\Big[E^{*}_\pm-\sqrt{k^{2}+m_{\mathrm{bare}}^{2}}\Big]\\
&-\frac{\big[\tfrac12(m^{*}_++m^{*}_-)-m_{\mathrm{bare}}\big]^{2}}{2\,\GStot}-\frac{\big[\tfrac12(m^{*}_+-m^{*}_-)\big]^{2}}{2\,\GSdif} .
\end{aligned}
\end{equation}
For the sets A and B the more attractive of the two scalar couplings is $\GS^{\mathrm{HF}}$ in both
forms ($\GStot$ in form Q, $\GSdif$ in form C), so that the critical strength
$g_{\mathrm{crit}}=G_c\,g/\GS^{\mathrm{HF}}$ of a field without mass term is the same for both
(Table~\ref{tab:vachalves}). Above it, \emph{in form Q the two halves acquire the same mass, a mass
term proportional to $\bar\psi\psi$, the mass term of the positive Lagrangian; in form C they acquire
masses of opposite sign, by Eq.~(\ref{eq:covmass}) a mass term proportional to $\tilde\psi\psi$, the
mass term of the covariant Lagrangian.} Each form generates dynamically the mass term that belongs
to its own adjoint; the magnitude of the mass is the same. With a bare mass, whose term is
$-m_{\mathrm{bare}}\bar\psi\psi$ in both forms, the two forms differ in kind
(Table~\ref{tab:vachalvesclock}, Fig.~\ref{fig:vachalves}). In form Q the bare mass and the dynamical
mass have the same structure, and the mass grows smoothly through the critical strength. In form C
the two masses are equal and slightly smaller than the bare mass up to a definite strength, because
$\GStot>0$ repels the total scalar density, and there they separate: a sharp transition at which the
symmetry between the halves is broken spontaneously. With the bare mass $0.05\,\Lambda$ it lies at
{{GEN:fm_vac_split_list}}, counted from the critical strength of the massless field, and agrees with
the linear criterion $1+2\,\GSdif\cdot4\int_{|\vec k|<\Lambda}d^{D}k\,(k^{2}/E^{*3})/(2\pi)^{D}=0$ at
the common mass to {{GEN:fm_vac_split_dev}}.

\begin{table}
\caption{\label{tab:vachalves}The vacuum of a field without mass term in the two forms, with the
filled states of negative energy (units of the cutoff): the critical strength of the sets A and B
[for the critical coupling $G_c$ of Eq.~(\ref{eq:Gc}) with $\nu=8$], and the masses
$(m^{*}_+,m^{*}_-)$ of the two halves at a strength 2 per cent above it (the representative with
$m^{*}_+>0$ of the degenerate solutions).}
\begin{ruledtabular}
\begin{tabular}{ccccc}
$D$ & set & $g_{\mathrm{crit}}\Lambda^{D-1}$ & form Q & form C \\
\hline
{{ROWS:tbl_fm_vacuum_critical|cols=0,1,3,4,5}}
\end{tabular}
\end{ruledtabular}
\end{table}

\begin{table*}
\caption{\label{tab:vachalvesclock}The masses of the two halves in vacuum, in units of the cutoff,
along the deflation clock counted from the critical strength of the massless field (set A), without
bare mass and with the bare mass $0.05\,\Lambda$. A single number: both halves have this mass; a
pair: the two masses, the larger first (which half has which is free by the symmetry between the
halves).}
\begin{ruledtabular}
\begin{tabular}{ccccccc}
$D$ & $m_{\mathrm{bare}}/\Lambda$ & form & $\td=-0.5$ & $0.1$ & $0.5$ & $1$ \\
\hline
{{ROWS:tbl_fm_vacuum}}
\end{tabular}
\end{ruledtabular}
\end{table*}

\begin{figure*}
\includegraphics[width=0.86\textwidth]{fig_forms_vacuum}
\caption{\label{fig:vachalves}The masses of the particles of the two halves in vacuum with the filled
states of negative energy, along the clock counted from the critical strength of a massless field
(dotted line; $D=4$, set A; units of the cutoff). (a) A field without mass term: above the critical
strength form Q gives both halves the same mass, form C gives them masses of opposite sign. (b) With
the bare mass $0.05\,\Lambda$: form Q passes smoothly through the critical strength; in form C the
two masses are equal up to a sharp transition and separate there.}
\end{figure*}

\emph{The thresholds of Sec.~\ref{sec:ks} and the vacuum.} The calculation of Sec.~\ref{sec:ks}
without the filled states is consistent only while the vacuum is stable, that is, while the most
attractive scalar coupling is weaker than $G_c$ for the cutoff of the model. The last column of
Table~\ref{tab:ksthr} gives the largest cutoff for which this holds at the binding threshold. It is
below the mass in all eight cases, and smaller in form C than in form Q, because form C needs
stronger couplings to bind matter: \emph{in both forms, and more so in form C, self-bound matter by
contact couplings alone presupposes a cutoff below the mass of the particles.} How the filled states
change a saturation point when they are included was computed for an example of one kind of
particle on the $7+1$ spacetime ($D=7$), with the bare mass adjusted so that the mass in vacuum
stays the same: the saturation point moves little while the cutoff stays below about half the mass
in vacuum, and the binding grows for larger cutoffs [Fig.~\ref{fig:cjtsea}(b)]. The results of
Sec.~\ref{sec:ks} are therefore those of an effective theory with a low cutoff.
%%</1644>

\subsection{Bare mass, self-interaction and Vorton-Scalar together}
\label{sec:unified}

Now let all terms of the model (\ref{eq:LVS}) act together, for fields in ordinary space ($D=3$) and
without the filled states. Let $\phi$ be the expectation value of the Vorton-Scalar,
$\langle\PhiVS\rangle=\phi/\sqrt2$, and $x=\phi/\vVS$. The mass of a fermion in a region with the
expectation value $\phi$ is $m_D+m_Y\,x$; let $f_D=m_D/m$ be the fraction of the mass in vacuum that
is a bare Dirac mass. In units $m=1$ the energy density of matter of the Fermi momentum $\kF$ is the
stationary value of
\begin{equation}
\begin{aligned}
\varepsilon={}&\frac{(1-f_D)^{2}\,(x^{2}-1)^{2}}{8\,C}+\langle E^{*}\rangle+\tfrac12\GV\,n^{2}\\
&-\frac{\big[\mst-f_D-(1-f_D)\,x\big]^{2}}{2\,\GS}
\end{aligned}
\label{eq:unified}
\end{equation}
with respect to $x$ and $\mst$: the first term is the potential of the scalar, with the number $C$
of Eq.~(\ref{eq:DC}); the others are Eq.~(\ref{eq:epsKS}) with the mass $f_D+(1-f_D)x$ in place of
$m$.\SIG{ Equation~(\ref{eq:unified}) is written for one kind of particle; for the two halves of the
field the last three terms are replaced by the energy density (\ref{eq:epshalves}) with
$f_D+(1-f_D)x$ in place of $m$, because the Yukawa coupling shifts the masses of both halves alike
(Sec.~\ref{sec:clock}).}{} For $f_D=0$ and $\GS=\GV=0$ this is the functional of Sec.~\ref{sec:cjt}, and the program
reproduces its masses and energies to {{GEN:uni_u1_dev}} and its binding threshold
$C=3\pi^{2}/8$.

\emph{At low density the Vorton-Scalar acts as a scalar contact coupling.} Expanding
Eq.~(\ref{eq:unified}) for small densities gives $\mst=m+(\GS-C/m^{2})\,\nS$: the exchange of the
radial mode adds $-\gVS^{2}/m_s^{2}$ to the scalar coupling, whatever fraction of the mass is a bare
mass (verified for five sets of parameters to {{GEN:uni_u2_dev}} at Fermi momenta up to a
twentieth of $m$). The scalar
channel of the self-interaction and the exchange of the radial mode are the same interaction at
momenta far below $m_s$; they differ at higher densities, where the scalar field itself changes.\SIG{ For
the two halves this means that the Vorton-Scalar adds $-C/m^{2}$ to the coupling $\GStot$ of the
total scalar density, in both forms.}{}

%%<1671>
\emph{The moment at which matter becomes self-bound.} All three couplings $C$, $\GS$, $\GV$ grow as
$e^{\td}$, so that their ratios are fixed and the deflation carries matter through a sequence of
states with growing $C$. Table~\ref{tab:unified} gives the smallest $C$ at which matter is
self-bound, and the shift of that moment relative to the model with the Vorton-Scalar alone
[Fig.~\ref{fig:unified}(a,b)]. \emph{A bare Dirac mass makes matter bind earlier}, by
{{GEN:uni_dtau_fd025}} units of the clock if a quarter of the mass is a bare mass and by
{{GEN:uni_dtau_fd05}} if half of it is: the matter no longer has to pay for the destruction of the
whole expectation value of the scalar. \emph{A vector repulsion delays the binding}, by
{{GEN:uni_dtau_kv01}} and {{GEN:uni_dtau_kv03}} units for the two vector couplings of the table,
without bare mass. A scalar contact attraction acts only where the particles keep a mass: without a
bare mass the matter of lowest energy at the threshold consists of massless particles, whose scalar
density vanishes, and the scalar coupling of the second row of the table changes nothing.

\begin{table*}
\caption{\label{tab:unified}Bare Dirac mass, contact couplings and Vorton-Scalar together ($D=3$,
without the filled states): the smallest $C$ at which uniform matter is self-bound, for three
vector couplings, and in parentheses the shift of the moment of binding along the deflation clock
relative to the first entry (negative: earlier).}
\begin{ruledtabular}
\begin{tabular}{ccccc}
$f_D$ & $-\GS/C$ & $\GV=0$ & $\GV=0.1\,C$ & $\GV=0.3\,C$ \\
\hline
{{ROWS:tbl_unified}}
\end{tabular}
\end{ruledtabular}
\end{table*}
%%</1671>
%%<1644>
\emph{The moment at which matter becomes self-bound.} All couplings grow as $e^{\td}$, so that their
ratios are fixed and the deflation carries matter through a sequence of states with growing $C$.
Table~\ref{tab:unified} gives the smallest $C$ at which matter is self-bound in the two forms, for
contact couplings $g_0<0$ and $g_1>0$ of the Lagrangian proportional to $C$, and the shift of that
moment relative to the model with the Vorton-Scalar alone. \emph{A bare Dirac mass makes matter bind
earlier}, by {{GEN:uni_dtau_fd025}} units of the clock if a quarter of the mass is a bare mass and by
{{GEN:uni_dtau_fd05}} if half of it is: the matter no longer has to pay for the destruction of the
whole expectation value of the scalar. In all {{GEN:fm_uni_n}} cases of the table the bound state at
the threshold has equal halves, in both forms: the attraction that binds it is that of the scalar
field, which acts on both halves alike, and contact couplings of these strengths are too weak to
polarize it. The threshold is that of matter of one kind with the couplings $\GStot$, $\GVtot$ of the
form, confirmed with the functional of two halves minimized over the density, the field and the
polarization. The forms differ because the couplings of the total densities differ
(Table~\ref{tab:kscoup}). In form Q the scalar contact coupling adds its attraction
$\GS^{\mathrm{HF}}$ to that of the scalar field, and a vector coupling adds the repulsion
$\GV^{\mathrm{HF}}$ [Fig.~\ref{fig:unified}(a,b), drawn against the mean-field couplings of matter of
one kind]. In form C the couplings of the total densities are $\GS^{\mathrm c}$ and
$\GV^{\mathrm c}$, both small, so that the contact couplings change the threshold less; which form
binds earlier depends on the case. Without bare mass the bound state of form C at the threshold
consists of massless particles ($x=0$), whose scalar densities vanish, and the threshold does not
depend on the vector coupling of the table.

\begin{table*}
\caption{\label{tab:unified}Bare Dirac mass, contact couplings and Vorton-Scalar together in the two
forms ($D=3$, without the filled states): the smallest $C$ at which uniform matter is self-bound,
in parentheses the shift of the moment of binding along the deflation clock relative to the first
row (negative: earlier), and the value $x$ of the scalar field in the bound state just above the
threshold.}
\begin{ruledtabular}
\begin{tabular}{ccccccc}
$f_D$ & $-g_0/C$ & $g_1/C$ & form C: $C$ ($\Delta\td$) & $x$ & form Q: $C$ ($\Delta\td$) & $x$ \\
\hline
{{ROWS:tbl_fm_unified}}
\end{tabular}
\end{ruledtabular}
\end{table*}
%%</1644>

\begin{figure*}
\includegraphics[width=\textwidth]{fig_unified}
\caption{\label{fig:unified}Bare Dirac mass, contact couplings and Vorton-Scalar together. (a), (b)
The shift of the moment at which matter becomes self-bound, in units of the deflation clock, against
the vector contact coupling, for three fractions $f_D$ of bare mass, without and with a scalar
contact coupling. (c) The square of the mass of the phase mode in matter divided by its value in
vacuum, against the Fermi momentum in units of the cutoff.}
\end{figure*}

\subsection{The mass of the phase mode}
\label{sec:phasemass}

A bare Dirac mass breaks the chiral symmetry $\UVS$ explicitly, and the phase mode of the
Vorton-Scalar is then no longer massless. Write $\PhiVS=(\vVS/\sqrt2)\,e^{i\vartheta_{\mathrm{VS}}}$
with a constant phase $\vartheta_{\mathrm{VS}}$. The mass term of the fermions becomes
$m_D+m_Y\,e^{i\vartheta_{\mathrm{VS}}\chi}$, the mass of a particle
\begin{equation}
m(\vartheta_{\mathrm{VS}})=\sqrt{m_D^{2}+m_Y^{2}+2\,m_Dm_Y\cos\vartheta_{\mathrm{VS}}},
\end{equation}
largest at $\vartheta_{\mathrm{VS}}=0$. The energy of the filled states of negative energy is lowest
where the mass is largest, so that the vacuum selects $\vartheta_{\mathrm{VS}}=0$, and the second
derivative of the energy density with respect to the phase is $\mth^{2}\vVS^{2}$. With
$dm/d\vartheta_{\mathrm{VS}}=0$ and $d^{2}m/d\vartheta_{\mathrm{VS}}^{2}=-m_Dm_Y/m$ at
$\vartheta_{\mathrm{VS}}=0$, and with the derivative of the energy density with respect to the mass,
which is the scalar density,
\begin{equation}
\mth^{2}\,\vVS^{2}=\Big(\big|\langle\bar\psi\psi\rangle_{\mathrm{vac}}\big|-\nS\Big)\,\frac{m_D\,m_Y}{m} .
\label{eq:mtheta}
\end{equation}
In vacuum ($\nS=0$) this is the relation of Gell-Mann, Oakes and Renner \cite{GellMannOakesRenner1968}
for the model at hand: the square of the mass of the phase mode is proportional to the bare mass and
to the scalar density of the vacuum. In ordinary space
$|\langle\bar\psi\psi\rangle_{\mathrm{vac}}|=(\nu m/4\pi^{2})\,[\Lambda\sqrt{\Lambda^{2}+m^{2}}-m^{2}\,\mathrm{arsinh}(\Lambda/m)]$;
the program confirms Eq.~(\ref{eq:mtheta}) against the numerical second derivative of the energy in
{{GEN:uni_u3_cases}} cases to {{GEN:uni_u3_dev}} of the value in vacuum. Three consequences follow.

\emph{The phase mode becomes heavier along the clock.} In the couplings of the model,
$\mth^{2}=C\,m_s^{2}\,f_D\,(|\langle\bar\psi\psi\rangle_{\mathrm{vac}}|-\nS)/[(1-f_D)\,m^{3}]$, which
grows as $e^{\td}$: $\mth\propto e^{\td/2}$, while all other masses stay fixed. Its value depends on
the cutoff through the scalar density of the vacuum (Table~\ref{tab:phasemode}).

\emph{In matter the phase mode is lighter.} The matter subtracts its scalar density:
$\mth^{2}(\kF)/\mth^{2}(0)=1-\nS/|\langle\bar\psi\psi\rangle_{\mathrm{vac}}|$
[Table~\ref{tab:phasemode}, Fig.~\ref{fig:unified}(c)].

\emph{A massive phase mode binds pairs less.} The exchange of the phase mode contributes to the
pairing potential (\ref{eq:VVS}) with the range $1/\mth$; a mass of three tenths of the fermion
mass reduces the gap of the example of Table~\ref{tab:gapcontact} ($m_s=m$) from
{{GEN:gapc_exchange}} to {{GEN:gapc_phase}} of the Fermi kinetic energy.

\begin{table}
\caption{\label{tab:phasemode}The phase mode of the Vorton-Scalar when a tenth of the fermion mass
is a bare Dirac mass ($f_D=0.1$; $D=3$; units $m=1$): scalar density of the vacuum, square of the
mass of the phase mode in vacuum in units of $C\,m_s^{2}$, and its reduction in matter of the Fermi
momenta $\Lambda/4$, $\Lambda/2$ and $3\Lambda/4$.}
\begin{ruledtabular}
\begin{tabular}{cccccc}
$\Lambda$ & $|\langle\bar\psi\psi\rangle_{\mathrm{vac}}|$ & $\mth^{2}/(C\,m_s^{2})$ & $\kF=\Lambda/4$ & $\Lambda/2$ & $3\Lambda/4$ \\
\hline
{{ROWS:tbl_phase_mode}}
\end{tabular}
\end{ruledtabular}
\end{table}
